\section{Introduction}
\label{sec:introduction}

Generative systems such as Suno and Udio synthesise complete songs---vocals,
instrumentation, lyrics and arrangement---that listeners struggle to
distinguish from professionally produced human music. Streaming platforms now
report that a substantial fraction of newly uploaded tracks are machine
generated, which turns detection from an academic question into an operational
requirement for distributors, rights organisations and digital service
providers.

Detection is also a forensic problem, and it inherits the asymmetry of one. A
missed synthetic track costs a platform little. A human artist wrongly flagged
as synthetic may lose income or reputation on the strength of a model's
confidence. Any detector proposed for this setting therefore has to be
characterised not only by how well it separates classes in the conditions it was
trained on, but by how it fails when those conditions change, and by whether the
confidence it reports can be trusted at the operating point where the decision
is actually made.

The literature reports a single headline failure: cross-generator collapse.
Cros Vila et al.~\cite{crosvila2025} benchmark full-song detectors under an
explicit cross-generator protocol and observe $F_1$ falling from 0.99
in-distribution to 0.629 on an unseen generative family. That number is widely
cited as the measure of the generalisation problem. We argue that it conflates
two distinct failures, and that separating them changes what the problem is.

This paper makes four contributions.

\begin{enumerate}
  \item We build a detector on the frozen MERT music foundation
        model~\cite{li2024mert} with a small attention-pooling head, and
        evaluate it under a protocol that varies the generator and the corpus of
        \emph{real} music independently. An ablation shows the head's design is
        not what produces its accuracy: discarding twelve of the thirteen hidden
        states, or replacing attention pooling with a mean, leaves every metric
        unchanged.
  \item We show that these two axes behave completely differently. Withholding
        an entire generator while holding the real class fixed costs almost
        nothing ($F_1$ degradation $-0.037$ and $-0.023$, against $-0.361$
        reported previously). Changing the corpus of real music collapses the
        same detector to near chance (AUROC 0.598) and causes it to label every
        genuine recording as synthetic (FPR 1.00).
  \item We give the explanation that reconciles them: the detector encodes what
        the real music in its training corpus looks like, rather than what
        generation leaves behind. Generalisation to new generators is therefore
        not evidence of generalisation at all.
  \item We report two diagnostics that any future work in this area can apply: a
        shortcut test showing that global loudness and dynamic-range descriptors
        alone separate the classes at 0.925 AUROC before normalisation, and a
        false-positive-rate check on unseen real music that exposes the
        corpus-dependence directly.
\end{enumerate}

We also find that the $F_1$ conventionally reported in this setting can be
actively misleading: on our cross-corpus evaluation the detector's $F_1$ of
0.750 is identical, to four decimal places, to that of a classifier that labels
every input synthetic. Reporting $F_1$ alone would have presented a total
failure as an improvement over the published baseline.
